\documentclass[11pt,a4paper]{article}
\usepackage[margin=2.2cm]{geometry}
\usepackage[T1]{fontenc}
\usepackage{lmodern,microtype,booktabs,tabularx,enumitem,xcolor,titlesec,fancyhdr}
\usepackage[hidelinks]{hyperref}
\setlength{\parindent}{0pt}\setlength{\parskip}{5pt}
\linespread{1.08}
\definecolor{ac}{HTML}{1F4E79}
\titleformat{\section}{\Large\bfseries\color{ac}}{\thesection}{0.6em}{}
\titleformat{\subsection}{\large\bfseries}{\thesubsection}{0.6em}{}
\setlist{itemsep=1pt,topsep=2pt}
\newcolumntype{L}{>{\raggedright\arraybackslash}X}
\pagestyle{fancy}\fancyhf{}\fancyfoot[C]{\thepage}\fancyhead[L]{\small Klaro study notes}\renewcommand{\headrulewidth}{0.3pt}
\newcommand{\key}[1]{\par\vspace{2pt}\noindent\fcolorbox{ac}{ac!8}{\parbox{\dimexpr\linewidth-2\fboxsep-2\fboxrule}{\textbf{Remember:} #1}}\par}
\newcommand{\qa}[2]{\par\vspace{3pt}\noindent\textbf{\color{ac}Q. #1}\par\noindent #2\par}

\begin{document}
{\LARGE\bfseries\color{ac} Papers behind Klaro: short summaries}\par
\vspace{2pt}Study guide 2 of 3. The report has 19 references: 12 papers and 7 tool pages.

\textbf{Read this first.} We did not publish a research paper ourselves. These are the papers \emph{cited} in Chapter 2 of the report (Background Study). For each one: what it did, what it tells us, and exactly how Klaro uses it. The first part of each entry matches what the report says. Extra detail marked $\dagger$ comes from general knowledge of the paper, not from the report, so open the paper if a teacher asks for specifics.

\key{Chapter 2 story: (A) people already analysed tenders, but only English and mostly with word counting; (B) long documents are hard for transformers; (C) small models make things up; (D) offline translation exists. Nobody had combined offline translation + small summary model + page-level checking on a laptop. That is the gap Klaro fills.}

\section{A. Earlier work on tender documents}

\subsection*{1. Beason et al., 2021 (SMU Data Science Review)}
\textbf{Automated analysis of RFPs using NLP for the technology domain.}\\
Processed a large set of U.S.\ Department of Defense and Homeland Security RFPs to spot technology trends.\\
\textbf{Tells us:} RFPs have no common layout and are noisy (contents pages, boilerplate, legal text), so cleaning comes first. Readers prefer exact structured facts over a free-flowing summary.\\
\textbf{Klaro uses:} skips contents pages; fixed nine headings.

\subsection*{2. Trappey et al., 2022 (J.\ Global Information Management)}
\textbf{Intelligent RFQ summarization using NLP, text mining and machine learning.}\\
Summarised more than 1,300 power-transformer tenders. Found key terms with \textbf{TF-IDF}, grouped sentences by topic, picked top sentences with \textbf{TextRank}.\\
\textbf{Tells us:} keeping the top 35\% of sentences kept about 80\% of important requirements and cut reading time to about 29\%.\\
\textbf{Klaro uses:} summarise topic by topic (our nine headings, two question sets per chunk) and aim for a real cut in length. (This is also why the context-size bug mattered: our summaries were not getting shorter.)\\
\textbf{Difference:} theirs is extractive (picks sentences); Klaro uses a language model to write.

\subsection*{3. Knisely and Pavliscsak, 2023 (Scientometrics)}
\textbf{Research proposal content extraction using NLP and semi-supervised clustering.}\\
Sorted research proposals into themes. A few proposals labelled by experts guided the sorting of the rest.\\
\textbf{Tells us:} a little human labelling beats none at all.\\
\textbf{Klaro uses:} this is the reason we built the golden reference set (and why we say it still needs human verification).

\section{B. Reading very long documents}

\subsection*{4. Vaswani et al., 2017 (NeurIPS): ``Attention is all you need''}
Introduced the \textbf{Transformer}, the design behind nearly all modern language models, built on self-attention.\\
\textbf{Tells us:} attention compares every token with every other token, so cost grows quickly (quadratic$^\dagger$) as text gets longer.\\
\textbf{Klaro uses:} the reason to cut an RFP into chunks of about 1,800 words instead of reading it whole.

\subsection*{5. Beltagy, Peters, Cohan, 2020: Longformer}
Changes attention so cost grows roughly linearly: each token looks at a local window plus a few global tokens$^\dagger$.\\
\textbf{Tells us:} thousands of words can be read in one pass.\\
\textbf{Klaro:} we cannot use it. A small local model like Qwen3.5 on a laptop cannot read an 84-page RFP in one go, so we chunk and merge.

\subsection*{6. Douzon et al., 2023 (ICDAR workshops)}
\textbf{Long-range transformer architectures for document understanding.}\\
Used the page layout of business documents, so a table value stays linked to its heading.\\
\textbf{Klaro uses:} the idea of keeping structure and page position. We keep page numbers on every chunk and fact. (We lose table layout, which is a stated limitation.)

\section{C. Small models and made-up answers}

\subsection*{7. Ji et al., 2023 (ACM Computing Surveys)}
\textbf{Survey of hallucination in natural language generation.}\\
A survey of why and how text generators state things that are not in the source.$^\dagger$ It separates statements that contradict the source from statements that cannot be checked against it.\\
\textbf{Klaro uses:} the definition of ``made-up information'' in our nomenclature and our hallucination check.

\subsection*{8. Maynez et al., 2020 (ACL)}
\textbf{On faithfulness and factuality in abstractive summarization.}\\
Found that abstractive summaries often contain content not supported by the document, and that common overlap scores like ROUGE hardly measure this$^\dagger$.\\
\textbf{Klaro uses:} the reason to check every point against the extracted facts and to be honest that our checks compare words only.

\subsection*{9. Liu et al., 2024 (TACL): ``Lost in the middle''}
\textbf{Tells us:} models use information at the start and end of a long input better than information in the middle$^\dagger$.\\
\textbf{Klaro uses:} another reason for chunks, and for our context-size finding: long inputs hurt in practice, not only in theory. (Our G3 run lost facts when given a very long fact list, which fits this idea, but we did not prove that is the cause.)

\subsection*{10. Lewis et al., 2020 (NeurIPS): Retrieval-augmented generation (RAG)}
First retrieve the relevant passages, then give only those to the generator.\\
\textbf{Klaro uses:} not RAG itself. Our remedy is different: force a page citation for every point and verify each point against collected facts. RAG remains an option for later (search over past proposals).

\subsection*{11. Hu et al., 2022 (ICLR): LoRA}
Freeze the big model and train small low-rank add-on matrices.\\
\textbf{Tells us:} far fewer trainable numbers and much less GPU memory than full fine-tuning (for GPT-3 the paper reports about 10,000 times fewer trainable parameters$^\dagger$).\\
\textbf{Klaro uses:} the fine-tuning kit. It is the only realistic route on an 8 GB card. Output is a small ``adapter''.

\section{D. Machine translation}

\subsection*{12. Tiedemann and Thottingal, 2020 (EAMT): OPUS-MT}
Open, small translation models, one per language direction, from Helsinki-NLP (built on the Marian toolkit$^\dagger$).\\
\textbf{Klaro uses:} \texttt{opus-mt-de-en} (about 301 MB) to translate German RFPs, and \texttt{opus-mt-en-de} for the German summary view. Small enough to run offline on a laptop.\\
\textbf{Note:} we cite it for what it is. We have not measured its quality ourselves and no German speaker has judged it.

\section{Tool references (not papers)}
\begin{tabularx}{\linewidth}{@{}p{2.6cm}L@{}}\toprule
Ollama & Runs open models locally and serves them on \texttt{127.0.0.1:11434}. Our model runner.\\
Qwen3.5 (4B) & The model, tag \texttt{qwen3.5:4b}, about 3.4 GB.\\
Electron & Desktop app framework (Chromium + Node). We follow its security guidance.\\
FastAPI & Python web framework for the local API.\\
Pydantic & Checks that model output has the right structure.\\
pypdf & Reads text from PDFs page by page.\\
LangGraph & Workflow library we \emph{chose not to use}: our steps are a fixed sequence, a few lines of Python run them.\\\bottomrule
\end{tabularx}

\section{One-line cheat sheet}
\begin{tabularx}{\linewidth}{@{}p{3.3cm}L@{}}\toprule
Beason 2021 & RFPs are messy; structured facts beat free summaries.\\
Trappey 2022 & TF-IDF + TextRank on 1,300 tenders; 35\% of sentences keeps about 80\% of requirements.\\
Knisely 2023 & A few expert labels improve clustering a lot.\\
Vaswani 2017 & Transformer; attention gets expensive with length.\\
Longformer 2020 & Cheaper attention for long text.\\
Douzon 2023 & Use document layout in long documents.\\
Ji 2023 & Survey of hallucination.\\
Maynez 2020 & Summaries often unfaithful; need a faithfulness check.\\
Liu 2024 & Models lose the middle of long inputs.\\
Lewis 2020 & RAG: retrieve then generate.\\
Hu 2022 & LoRA: train a small add-on, not the whole model.\\
Tiedemann 2020 & OPUS-MT: small open translation models.\\\bottomrule
\end{tabularx}
\end{document}
